\documentclass[10pt]{article}
\usepackage[margin=0.85in]{geometry}
\usepackage[T1]{fontenc}
\usepackage{lmodern,microtype,booktabs,longtable,array}
\usepackage[round,authoryear]{natbib}
\usepackage[hidelinks]{hyperref}
\newcommand{\PaperTitle}{Geographic variation in local-government reservation assignment}

\title{Supporting information for\\\PaperTitle}
\author{Gaurav Sood}
\date{\today}
\renewcommand{\thesection}{S\arabic{section}}
\renewcommand{\thepage}{S\arabic{page}}
\begin{document}
\maketitle
This supporting information documents the assignment claims examined in the research note. It screens the upstream literature corpus, its available versions and supplements, and assignment-specific sources traced from the research note. Each assertion is attached to its source, geography, period, office, and stated qualifications. An identifying assumption, a legal prescription, and a documented administrative practice are recorded separately. Missing historical inputs do not turn a claim into a refuted claim.

The machine-readable source records retain original paths, URLs and SHA-256 hashes. PDFs remain upstream. Short quotations preserve exact wording; the accompanying paraphrases recover the procedure and qualifications. Claims marked outside current coverage concern a different population or period. Printed and PDF page numbers are distinguished.
Some locally supplied drafts lack a verified public URL or request noncirculation. Their status is recorded in the screening notes. No quotations from those drafts are reproduced, and they do not enter the main note's comparison. They should not be treated as independently accessible published evidence.
\section{Source screening}
\footnotesize\input{build/claims_screening.tex}\normalsize
\input{build/claims_records.tex}
\section{Claims ready for a comparison with current data}
\footnotesize\input{build/claims_next_tests.tex}\normalsize
\bibliographystyle{plainnat}
\bibliography{references,claims_a,claims_b,claims_external}
\end{document}
